%% ============================================================
%% FRONT MATTER
%% ============================================================
\pagestyle{plain}
\pagenumbering{roman}

%% ============================================================
%% ABSTRACT
%% ============================================================
\chapter*{Abstract}
\addcontentsline{toc}{chapter}{Abstract}

This thesis compares three power cycles for industrial waste heat recovery
(WHR) in an Australian context --- the Steam Rankine Cycle (SRC), the
Supercritical CO$_2$ (sCO$_2$) Brayton cycle and the Transcritical CO$_2$
(tCO$_2$) Rankine cycle --- under one consistent set of thermodynamic,
sizing, cost and economic assumptions, applied to three representative
sites: LNG gas processing in the Pilbara (200\,\textdegree C, 30\,kg/s),
cement preheater exhaust in Gladstone (350\,\textdegree C, 50\,kg/s) and
electric arc furnace off-gas in Sydney (500\,\textdegree C, 20\,kg/s). Four
research questions address thermodynamic performance (RQ1), equipment sizing
and capital cost (RQ2), techno-economic viability (RQ3) and off-design
flexibility (RQ4). All models are implemented in Python with CoolProp and
every result is regenerated by a single script.

Thirty findings are reported. With the recuperator energy balance limited by
both streams --- a second-law requirement that the interim analysis did not
satisfy --- the steam cycle has the highest thermal efficiency at every
turbine inlet temperature from 150 to 600\,\textdegree C (19.1--35.4\,\%
against 13.3--30.1\,\% for tCO$_2$ and 5.3--30.1\,\% for sCO$_2$), and no
crossover occurs within the WHR range. A 200-node integration of the tCO$_2$
heater against a waste-gas stream whose temperature is linear in heat duty
shows a double pinch at both terminals with the driving force widest in the
middle (42--65\,K), so that two-point LMTD sizing overestimates heater area
by 129--223\,\% when the pseudo-critical $c_p$ peak lies inside the heater.
The CO$_2$ cycles' decisive weakness in Australia is the heat sink: the
tCO$_2$ baseline condenses at 14.3\,\textdegree C, which a chiller could
supply only by consuming 63--134\,\% of the net output, and the sCO$_2$ cycle
loses 42\,\% of its 200\,\textdegree C efficiency if its compressor inlet
warms from 32 to 40\,\textdegree C.

Integrated with each site's stream through a 10\,K-approach heater, the
steam boiler must run at 0.16--3.7\,MPa rather than 10\,MPa, and the tCO$_2$
cycle recovers the most power at 200 and 350\,\textdegree C (723 and
3418\,kW against 412 and 2639\,kW for steam) because its cold heater inlet
cools the gas furthest; the recuperated sCO$_2$ cycle recovers the least at
every site. CO$_2$ turbines are 15--26 times more compact than the steam
turbine, but at equal output steam is the cheapest plant at every site
(AUD\,2{,}600--4{,}400/kW at 1\,MWe in 2026), tCO$_2$ costs 30--65\,\% more
and sCO$_2$ two to three times more. On a fixed-source basis steam gives the
lowest LCOE at all three sites (AUD\,96, 35 and 31/MWh) and the highest NPV
with 3--4-year discounted paybacks and 99.8--100\,\% Monte Carlo viability at
350 and 500\,\textdegree C; tCO$_2$ is the better investment at
200\,\textdegree C (NPV AUD\,1.71\,M against 1.12\,M) only where a natural
cold sink exists; sCO$_2$ is not viable at any site. The tCO$_2$ cycle is
the most flexible at part load (Marchionni score 0.78 against 0.73 for steam
and 0.71 for sCO$_2$).

The thesis concludes that the figure of merit for WHR cycle selection is the
net work recovered per kilogram of waste gas and its cost, not thermal
efficiency, and that at Australian sink temperatures the steam Rankine cycle
is the default choice above 300\,\textdegree C, with transcritical CO$_2$
reserved for low-temperature sources where cold water is available.

\bigskip\noindent\textbf{Keywords:} CO$_2$ power cycle; waste heat recovery;
supercritical CO$_2$; transcritical CO$_2$; steam Rankine cycle; exergy
analysis; pseudo-critical; pinch analysis; techno-economic analysis; LCOE;
Australia; energy transition.

\clearpage

%% ============================================================
%% EXECUTIVE SUMMARY
%% ============================================================
\chapter*{Executive Summary}
\addcontentsline{toc}{chapter}{Executive Summary}

This report evaluates the commercial viability of retrofitting waste heat
recovery (WHR) power cycles at three Australian industrial sites: a Pilbara
LNG train (200\,\textdegree C), a Gladstone cement preheater
(350\,\textdegree C) and a Sydney electric arc furnace (500\,\textdegree C).
Three cycles --- Steam Rankine (SRC), Supercritical CO$_2$ Brayton (sCO$_2$)
and Transcritical CO$_2$ Rankine (tCO$_2$) --- were modelled on one
consistent basis, each sized to what it can recover from the site's exhaust
stream and costed at that size in 2026 Australian dollars.

\textbf{Key findings for capital allocation:}
\begin{itemize}[leftmargin=1.5em,itemsep=6pt]
\item \textbf{High- and medium-temperature WHR (350--500\,\textdegree C):}
The steam Rankine cycle is the clear choice. It is the most efficient of the
three under Australian sink conditions, the cheapest to build
(AUD\,1{,}900--2{,}200/kW at the 2.3--2.6\,MWe site scale), and the most
bankable: LCOE of AUD\,31--35/MWh against site prices of 90--100/MWh,
discounted payback of 3.2--4.2 years, IRR of 28--36\,\%, and positive NPV in
99.8--100\,\% of 10,000 cost, price and utilisation scenarios.
\item \textbf{Low-temperature WHR (200\,\textdegree C):} Marginal for any
cycle. Steam and tCO$_2$ both return about 12\,\% IRR with a one-in-three
chance of loss, and depend on the Pilbara price premium. The tCO$_2$ cycle
recovers 75\,\% more electricity from the stream and has 53\,\% more NPV, but
only if a heat sink near 4\,\textdegree C is available; supplying it with a
chiller would consume more power than the plant makes.
\item \textbf{Supercritical CO$_2$:} Not investable in its simple recuperated
form at any of the three sites. With cooling water that cannot bring the
compressor inlet below 36\,\textdegree C, the cycle's back-work ratio is
31--64\,\%, its turbomachinery costs more than a whole steam plant, and its
recuperated heater leaves most of the exhaust heat unrecovered. NPV is
negative at two sites and marginal at the third.
\item \textbf{Flexibility:} Where the heat source fluctuates (batch
furnaces, intermittent trains), the tCO$_2$ cycle retains 56--67\,\% of its
efficiency at one-fifth load against 38--66\,\% for steam, and starts up most
simply.
\end{itemize}

\textbf{Investment recommendation:}\\
Target \textbf{steam Rankine} technology for Australian WHR sources at
350\,\textdegree C and above. Treat 200\,\textdegree C sources as
site-specific studies, with transcritical CO$_2$ preferred only where cold
water or a co-located refrigeration load exists. Do not pursue simple
supercritical CO$_2$ Brayton systems below 500\,\textdegree C in Australian
conditions. Evaluate every option on the electricity recovered from the
actual stream and its cost, not on thermal efficiency.

\clearpage

%% ============================================================
%% DECLARATION
%% ============================================================
\chapter*{Declaration by Author}
\addcontentsline{toc}{chapter}{Declaration by Author}
%% STUDENT: read every sentence of this declaration and edit it so that it
%% describes exactly what you did and what the AI tools did. Insert the model
%% name(s) exactly as shown in your Claude app. Sign only when it is true.
The work presented in this Final Thesis is my own except where I have stated
otherwise. The research questions, the choice of cycles, sites and modelling
assumptions, the interpretation of all results, and the engineering
conclusions and recommendations are mine, carried out under the supervision
of Dr Alexander Klimenko at the University of Queensland.

In completing this thesis I used two generative artificial intelligence (AI)
tools. Their use is disclosed in full here and in the AI Assistance log in
Appendix~\ref{app:ai}, in accordance with UQ's policy on the use of AI in
assessment~\cite{UQ2026AI}.

\bigskip
\noindent\textbf{Tool 1 --- Claude (Anthropic, 2026)~\cite{Anthropic2026}}

\noindent I used Claude for the following purposes:

\begin{itemize}[noitemsep,topsep=3pt]
  \item Drafting initial paragraphs for the background and literature review
    (Chapters~1 and~2), which I checked against the primary sources and
    revised.
  \item \LaTeX{} typesetting of the document, equations, tables and figure
    environments, and the Matplotlib styling of the figures.
  \item Reviewing the Python models I wrote for the Interim Report. This
    review identified three errors in my code: a recuperator energy balance
    that violated the second law, a heater gas-temperature profile that was
    linear in node index rather than in heat duty, and heater
    temperature-difference expressions in the cost model that produced
    unphysical values. It also identified that my turbine cost coefficients
    could not be traced to the source I had cited.
  \item Restructuring my model scripts into the \texttt{co2whr} package
    described in Appendix~\ref{app:code}, implementing the corrections above,
    adding the fixed-source site integration, the consistent cost basis of
    Table~\ref{tab:cost_basis}, discounted payback and the breakeven analysis.
    I ran the resulting code, checked the energy and exergy balances, traced
    every correlation to its source, and compared the outputs with my
    original results and with the literature.
  \item Drafting the text of Chapters~4--9, the abstract and the executive
    summary from the model outputs, which I then reviewed, corrected where
    necessary and approved. The findings, their numbering and their
    interpretation were settled by me.
\end{itemize}

\noindent The cycle models of the Interim Report (steam, sCO$_2$ and tCO$_2$
state-point models, exergy balances, pinch analysis, SP/VFR sizing, LCOE/NPV
and Monte Carlo, Stodola part-load model) were written by me and remain the
physical basis of the corrected package. Every number in this thesis was
produced by running that package, and I can reproduce and explain every
result and every line of the code independently of any AI tool.

\bigskip
\noindent\textbf{Tool 2 --- GitHub Copilot (Microsoft, 2026)~\cite{GitHubCopilot2026}}

\noindent I used GitHub Copilot for autocomplete assistance with Python
syntax while writing the original model files --- specifically CoolProp
\texttt{PropsSI} argument strings and NumPy and Matplotlib syntax. GitHub
Copilot did not generate any model logic, result or text appearing in this
thesis.

\bigskip
\noindent I have read and understood the University of Queensland's Student
Code of Conduct and the UQ policy on academic integrity and the use of AI in
assessment~\cite{UQ2026AI}. I confirm that this submission complies with
those policies and that I am able to demonstrate my learning, skills and
abilities regarding all content in this submission independently of any AI
tool.

\bigskip\bigskip\bigskip
\noindent\rule{7cm}{0.4pt}\\[0.2cm]
\textbf{Vennkkata Reddi Cheemalapati}\\
\textit{Student Number: 46555991}\\
\textit{Date: 22 October 2026}
\clearpage

%% ============================================================
%% ACKNOWLEDGEMENTS
%% ============================================================
\chapter*{Acknowledgements}
\addcontentsline{toc}{chapter}{Acknowledgements}
I am grateful to Dr Alexander Klimenko for his supervision throughout this
project. His direction to model transitional thermophysical effects in the
CO$_2$ heat exchanger, and his comments on the interim report and conference
presentation --- on the recuperator, on the economic basis and on the
Australian heat sink --- shaped the final analysis. I thank the School of
Mechanical and Mining Engineering for access to computational resources.

The open-source thermophysical property library CoolProp~\cite{Bell2014},
the Python scientific computing ecosystem (NumPy, SciPy, Matplotlib) and
NIST REFPROP data~\cite{Lemmon2018} were essential to the modelling work.

Two generative AI tools were used in preparing this thesis, in accordance
with UQ's policy on AI in assessment~\cite{UQ2026AI} and the UQ Library
guidance on acknowledging and referencing AI~\cite{UQLib2026}. Claude
(Anthropic, 2026)~\cite{Anthropic2026} assisted with drafting, typesetting,
code review and the restructuring of the model code, as set out in the
Declaration and in Appendix~\ref{app:ai}; GitHub Copilot (Microsoft,
2026)~\cite{GitHubCopilot2026} assisted with Python syntax. The research
design, the interpretation of results and the conclusions are my own.
\clearpage
